\documentclass[UTF8,a4paper,11pt]{ctexart}

\usepackage{amsmath,amssymb,mathtools}
\usepackage{geometry}
\usepackage{fancyhdr}
\usepackage{enumitem}
\usepackage{hyperref}

\geometry{left=2.5cm,right=2.5cm,top=2.5cm,bottom=2.5cm}
\setlength{\parindent}{2em}
\linespread{1.2}
\setlength{\headheight}{14pt}

\pagestyle{fancy}
\fancyhf{}
\fancyfoot[C]{\thepage}
\fancyhead[L]{数值分析\quad 第十次理论作业\quad 学号：241098117，姓名：张城宗}
\fancyhead[R]{第五章}

\newcommand{\dd}{\,\mathrm{d}}
\newcommand{\ee}{\mathrm{e}}
\newcommand{\QED}{\hfill$\square$}

\begin{document}

\begin{center}
  {\LARGE\bfseries 数值分析\quad 第十次理论作业}\\[6pt]
  {\large 第五章第 13、15、17、23、28、30 题}\\[4pt]
  {\normalsize 学号：241098117\quad 姓名：张城宗}
\end{center}

\bigskip

%=====================================================
\noindent\textbf{第五章\quad 第 13 题}\quad
试证明，解初值问题
\[
  y'=f(t,y),\qquad y(t_0)=\eta
\]
的隐式单步法
\[
  y_{n+1}=y_n+\frac{h}{6}\left[
  4f(t_n,y_n)+2f(t_{n+1},y_{n+1})+h f'(t_n,y_n)
  \right]
\]
为三阶方法。

\medskip
\noindent\textbf{证明}\quad
这里的 \(f'(t_n,y_n)\) 理解为沿微分方程解曲线的全导数，即
\[
  f'(t,y)=\frac{\dd}{\dd t}f(t,y(t))
  = f_t(t,y)+f_y(t,y)f(t,y).
\]
把精确解 \(y(t)\) 代入该格式。记
\[
  y_n=y(t_n),\quad
  y'_n=y'(t_n),\quad y''_n=y''(t_n),\quad \ldots .
\]
由微分方程有
\[
  f(t_n,y(t_n))=y'_n,\qquad
  f'(t_n,y(t_n))=y''_n,
\]
并且
\[
\begin{aligned}
  f(t_{n+1},y(t_{n+1}))
  &=y'(t_n+h)  \\
  &=y'_n+h y''_n+\frac{h^2}{2}y^{(3)}_n
    +\frac{h^3}{6}y^{(4)}_n+O(h^4).
\end{aligned}
\]
于是该格式右端给出的增量为
\[
\begin{aligned}
  \frac{h}{6}\left[
  4y'_n+2y'(t_n+h)+h y''_n
  \right]
  &=h y'_n+\frac{h^2}{2}y''_n
    +\frac{h^3}{6}y^{(3)}_n
    +\frac{h^4}{18}y^{(4)}_n+O(h^5).
\end{aligned}
\]
而精确解的 Taylor 展开为
\[
  y(t_n+h)-y(t_n)
  =h y'_n+\frac{h^2}{2}y''_n
    +\frac{h^3}{6}y^{(3)}_n
    +\frac{h^4}{24}y^{(4)}_n+O(h^5).
\]
二者相减得局部截断误差
\[
  \tau_{n+1}
  =-\frac{h^4}{72}y^{(4)}_n+O(h^5)=O(h^4).
\]
因此该单步法的阶数为 \(3\)。一般情形下 \(y^{(4)}_n\) 项不恒为零，所以它不是四阶方法。\QED

\bigskip\bigskip

%=====================================================
\noindent\textbf{第五章\quad 第 15 题}\quad
试写出经典的四阶 Runge--Kutta 方法解初值问题
\[
  y'=f(t),\qquad t_0\le t\le T,\qquad y(t_0)=y_0
\]
的计算公式。它与数值积分公式有什么关系？

\medskip
\noindent\textbf{解}\quad
经典四阶 Runge--Kutta 方法为
\[
\begin{aligned}
  K_1&=f(t_n),\\
  K_2&=f\!\left(t_n+\frac h2\right),\\
  K_3&=f\!\left(t_n+\frac h2\right),\\
  K_4&=f(t_n+h),\\
  y_{n+1}
  &=y_n+\frac h6(K_1+2K_2+2K_3+K_4).
\end{aligned}
\]
因为此处 \(f\) 与 \(y\) 无关，所以 \(K_2=K_3\)，从而
\[
  \boxed{
  y_{n+1}
  =y_n+\frac h6\left[
  f(t_n)+4f\!\left(t_n+\frac h2\right)+f(t_n+h)
  \right].}
\]
另一方面，精确解满足
\[
  y(t_{n+1})-y(t_n)=\int_{t_n}^{t_{n+1}}f(t)\dd t.
\]
上式正是用 Simpson 求积公式近似该积分。因此，当微分方程右端只依赖于 \(t\) 时，经典四阶 Runge--Kutta 方法退化为逐步应用 Simpson 公式的数值积分方法。

\bigskip\bigskip

%=====================================================
\noindent\textbf{第五章\quad 第 17 题}\quad
试证明经典四阶 Runge--Kutta 方法解初值问题
\[
  y'=\lambda y,\qquad y(t_0)=y_0
\]
的计算公式可写成
\[
  y_{n+1}=\left(
  1+\lambda h+\frac12(\lambda h)^2
  +\frac16(\lambda h)^3+\frac1{24}(\lambda h)^4
  \right)y_n,
\]
并对初值问题
\[
  y'=-10y,\qquad y(0)=1
\]
求 \(y(1)\) 的近似值，取步长 \(h=0.1\)。

\medskip
\noindent\textbf{解}\quad
令 \(z=\lambda h\)。经典四阶 Runge--Kutta 方法的四个斜率为
\[
\begin{aligned}
  K_1&=\lambda y_n,\\
  K_2&=\lambda\left(y_n+\frac h2K_1\right)
      =\lambda\left(1+\frac z2\right)y_n,\\
  K_3&=\lambda\left(y_n+\frac h2K_2\right)
      =\lambda\left(1+\frac z2+\frac{z^2}{4}\right)y_n,\\
  K_4&=\lambda\left(y_n+hK_3\right)
      =\lambda\left(1+z+\frac{z^2}{2}+\frac{z^3}{4}\right)y_n.
\end{aligned}
\]
代入
\[
  y_{n+1}=y_n+\frac h6(K_1+2K_2+2K_3+K_4)
\]
可得
\[
\begin{aligned}
  y_{n+1}
  &=\left[
  1+z+\frac{z^2}{2}+\frac{z^3}{6}+\frac{z^4}{24}
  \right]y_n \\
  &=\left(
  1+\lambda h+\frac12(\lambda h)^2
  +\frac16(\lambda h)^3+\frac1{24}(\lambda h)^4
  \right)y_n.
\end{aligned}
\]
这就是所要证明的公式。

对 \(y'=-10y\)，取 \(h=0.1\)，则 \(z=\lambda h=-1\)，所以每一步的放大因子为
\[
  R(-1)=1-1+\frac12-\frac16+\frac1{24}=\frac38.
\]
从 \(t=0\) 到 \(t=1\) 共 \(10\) 步，且 \(y_0=1\)，故
\[
  y(1)\approx y_{10}
  =\left(\frac38\right)^{10}
  =\frac{59049}{1073741824}
  \approx 5.49937\times 10^{-5}.
\]

\bigskip\bigskip

%=====================================================
\noindent\textbf{第五章\quad 第 23 题}\quad
应用 Heun 方法
\[
  y_{n+1}=y_n+\frac h4\left[
  f(t_n,y_n)
  +3f\left(t_n+\frac23h,\,
  y_n+\frac23h f(t_n,y_n)\right)
  \right]
\]
解初值问题
\[
  y'=-y,\qquad y(0)=y_0
\]
时，问步长 \(h\) 应取何值才能保证方法的绝对稳定性？

\medskip
\noindent\textbf{解}\quad
对试验方程 \(y'=-y\)，有 \(f(t,y)=-y\)。代入 Heun 方法，
\[
\begin{aligned}
  y_{n+1}
  &=y_n+\frac h4\left[
  -y_n
  +3\left(-y_n+\frac23h y_n\right)
  \right] \\
  &=\left(1-h+\frac{h^2}{2}\right)y_n.
\end{aligned}
\]
因此稳定放大因子为
\[
  R(h)=1-h+\frac{h^2}{2}.
\]
绝对稳定要求
\[
  |R(h)|<1.
\]
由于
\[
  R(h)=\frac12(h-1)^2+\frac12>0,
\]
所以只需
\[
  1-h+\frac{h^2}{2}<1.
\]
化简得
\[
  \frac{h^2}{2}-h<0
  \quad\Longleftrightarrow\quad
  h(h-2)<0.
\]
因此在 \(h>0\) 的步长条件下，
\[
  \boxed{0<h<2.}
\]

\bigskip\bigskip

%=====================================================
\noindent\textbf{第五章\quad 第 28 题}\quad
试从
\[
  y(t_{n+1})-y(t_{n-p})
  =\int_{t_{n-p}}^{t_{n+1}} f(t,y(t))\dd t
  =\int_{t_{n-p}}^{t_{n+1}} y'(t)\dd t
\]
导出解初值问题
\[
  y'=f(t,y),\qquad a\le t\le b,\qquad y(a)=y_0
\]
的 Nystrom 方法：
\[
  y_{n+1}=y_{n-1}+2hf_n.
\]

\medskip
\noindent\textbf{解}\quad
取 \(p=1\)，则
\[
  y(t_{n+1})-y(t_{n-1})
  =\int_{t_{n-1}}^{t_{n+1}} y'(t)\dd t.
\]
在区间 \([t_{n-1},t_{n+1}]\) 上用中点公式近似右端积分。该区间长度为 \(2h\)，中点为 \(t_n\)，所以
\[
  \int_{t_{n-1}}^{t_{n+1}} y'(t)\dd t
  \approx 2h\,y'(t_n)
  =2h\,f(t_n,y(t_n)).
\]
以数值解 \(y_n\) 近似 \(y(t_n)\)，并记
\[
  f_n=f(t_n,y_n),
\]
便得到
\[
  y_{n+1}-y_{n-1}=2h f_n,
\]
即
\[
  \boxed{y_{n+1}=y_{n-1}+2h f_n.}
\]
这就是 Nystrom 方法。

\bigskip\bigskip

%=====================================================
\noindent\textbf{第五章\quad 第 30 题}\quad
证明 BDF 方法
\[
  \sum_{m=1}^{k}\frac{\nabla^m y_n}{m}=h f(t_n,y_n)
  \tag{5.4.23}
\]
是 \(k\) 阶收敛的。

\medskip
\noindent\textbf{证明}\quad
BDF 方法的构造是：用经过
\[
  (t_n,y_n),(t_{n-1},y_{n-1}),\ldots,(t_{n-k},y_{n-k})
\]
的 \(k\) 次后差插值多项式 \(q_k(t)\) 的导数 \(q_k'(t_n)\) 来近似 \(y'(t_n)\)。令
\[
  s=\frac{t-t_n}{h}.
\]
Newton 后差插值多项式可写成
\[
  q_k(t_n+sh)
  =y_n+s\nabla y_n+\frac{s(s+1)}{2!}\nabla^2y_n
  +\cdots+
  \frac{s(s+1)\cdots(s+k-1)}{k!}\nabla^k y_n.
\]
对 \(t\) 求导，并令 \(s=0\)，注意
\[
  \left.\frac{\dd}{\dd s}
  \frac{s(s+1)\cdots(s+m-1)}{m!}\right|_{s=0}
  =\frac{(m-1)!}{m!}=\frac1m,
\]
于是
\[
  q_k'(t_n)
  =\frac1h\sum_{m=1}^{k}\frac{\nabla^m y_n}{m}.
\]
所以式 (5.4.23) 等价于
\[
  q_k'(t_n)=f(t_n,y_n).
\]

下面考察截断误差。把精确解 \(y(t)\) 代入，并令 \(q_k(t)\) 为它在
\[
  t_n,t_{n-1},\ldots,t_{n-k}
\]
上的 \(k\) 次插值多项式。插值导数误差满足
\[
  y'(t_n)-q_k'(t_n)=O(h^k).
\]
因此
\[
  \sum_{m=1}^{k}\frac{\nabla^m y(t_n)}{m}
  -h y'(t_n)
  =h\bigl(q_k'(t_n)-y'(t_n)\bigr)
  =O(h^{k+1}).
\]
这说明 BDF 方法的局部截断误差为 \(O(h^{k+1})\)，即方法具有 \(k\) 阶相容精度。

在线性多步法理论中，若出发值满足 \(O(h^k)\) 精度并且方法满足零稳定性，则局部截断误差 \(O(h^{k+1})\) 推出整体误差 \(O(h^k)\)。讲义中通常使用的 BDF 方法 \(k=1,\ldots,6\) 满足相应的零稳定条件，因此 BDF 方法 (5.4.23) 是 \(k\) 阶收敛的。\QED

\end{document}
